\section{Introduction}

Here is a confession: the first time I read the abstract of \emph{``Attention Is All You Need''} \cite{vaswani2017attention}, I understood maybe a third of it.
Not because I was stupid, but because the authors were writing for an audience of people who already knew what sequence transduction models were, who already had opinions about whether RNNs or CNNs were better for machine translation, and who would immediately understand what a BLEU score of 28.4 meant and why it was exciting.

I was not that audience.
If you are not that audience either, this paper is for you.

What I want to do is take the abstract --- just the abstract, all six sentences of it --- and unpack every single thing it says.
Not quickly, not superficially, but in a way where by the end you could explain each sentence to someone else and actually mean it.

For reference, here is what we are working with:

\begin{tcolorbox}[colback=gray!5!white, colframe=gray!60!black, title=The Original Abstract]
\small
The dominant sequence transduction models are based on complex recurrent or convolutional neural networks that include an encoder and a decoder. The best performing models also connect the encoder and decoder through an attention mechanism. We propose a new simple network architecture, the Transformer, based solely on attention mechanisms, dispensing with recurrence and convolutions entirely. Experiments on two machine translation tasks show these models to be superior in quality while being more parallelizable and requiring significantly less time to train. Our model achieves 28.4 BLEU on the WMT 2014 English-to-German translation task, improving over the existing best results, including ensembles, by over 2 BLEU. On the WMT 2014 English-to-French translation task, our model establishes a new single-model state-of-the-art BLEU score of 41.0 after training for 3.5 days on eight GPUs, a small fraction of the training costs of the best models from the literature.
\end{tcolorbox}

Six sentences.
Roughly a dozen concepts that a non-specialist would have no reason to know.
Let's go.
